\section{F-SM-43: the full original interval hypotheses do not force AB4}

\subsection{Statement and exact scope}
For every field \(k\), there exists a complete and cocomplete
\(k\)-linear abelian symmetric monoidal category \(\mathcal M\)
with enough projectives satisfying the scalar adjunction,
canonical unit-normalization, and full tensor-interval assumptions
of cote~134-8, footers~46--48, for which countable coproducts of
monomorphisms are not always monic.  The tensor is right exact
and preserves colimits in each variable; it is even closed.
The distinguished interval and its quotient by the unit are flat.
Thus none of these extra features forces AB4.

Nevertheless, the canonical derived scalar functor in this example
preserves arbitrary coproducts, its canonical counit is an
equivalence on every scalar module, and its ordinary localization
is \(D_{\ge0}(k)\).  This is a counterexample only to automatic
AB4, not to the original localization conclusion.

\subsection{Construction and complete proof}
Let \(\mathcal A\) be any complete and cocomplete \(k\)-linear
abelian category with enough projectives.  We shall at the end
take \(\mathcal A\) to be the explicit category of F-SM-34,
which fails countable AB4.  No tensor on \(\mathcal A\) is assumed.

For a vector space \(U\) and \(A\in\mathcal A\), let
\(U\odot A\) denote the canonical scalar copower, characterized by
\[
 \operatorname{Hom}_{\mathcal A}(U\odot A,B)
  =\operatorname{Hom}_k\bigl(U,
               \operatorname{Hom}_{\mathcal A}(A,B)\bigr).
                                                               \tag{SM43.1}
\]
It exists: after choosing a basis of \(U\), take the corresponding
coproduct of copies of \(A\).  The displayed universal property
makes the construction independent of that choice and functorial.
Currying in vector spaces and Yoneda give the canonical maps
\[
 (U\otimes V)\odot A\simeq U\odot(V\odot A),
       \qquad k\odot A\simeq A.                              \tag{SM43.2}
\]
Every iterated such map is the same iterated currying map under
(SM43.1), proving its associativity and unit coherence.
The action preserves colimits in each variable: in \(U\) this
follows from (SM43.1); for fixed \(U\), its basis description is
a coproduct functor on \(\mathcal A\), which preserves colimits.
No claim of exactness in \(A\) is made.

Set
\[
 \mathcal M=\operatorname{Ch}_{\ge0}(k)\times\mathcal A,
 \qquad
 (C,A)\boxtimes(D,B)=
  \bigl(C\otimes_k D, C_0\odot B\oplus D_0\odot A\bigr).
                                                               \tag{SM43.3}
\]
The product category is abelian, complete and cocomplete.
The tensor formula is a bifunctor, using the functorial copower
maps on its second component.  It is additive and preserves
colimits in each variable, since degree-zero evaluation of
connective chain complexes preserves colimits.  Hence it is
right exact in each variable.

Here are all its coherence maps.  The first component uses the
usual chain tensor, associator and Koszul symmetry.  The second
component of any parenthesization of an \(m\)-fold product has
the canonical normal form
\[
 \bigoplus_{i=1}^{m}
       \left(\bigotimes_{j\ne i}(C_j)_0\right)\odot A_i.
                                                               \tag{SM43.4}
\]
Define associators through this normal form, using (SM43.2).
Both pentagon composites are the same currying map on each
labelled summand.  Define the symmetry by permuting the summand
labels and the degree-zero scalar factors; its hexagon is the
equality of the corresponding permutations.  These maps are
natural by (SM43.1).  Their unit is
\(\mathbf1=(k[0],0)\), and the triangle reduces to
\(k\odot A=A\).  No Koszul sign enters the second component,
whose scalar factors all have degree zero.  This verifies the
full symmetric monoidal structure without a tensor on
\(\mathcal A\).

For completeness, closedness also follows directly.  The scalar
cotensor \(U\pitchfork B\) exists by products in \(\mathcal A\),
with
\[
 \operatorname{Hom}_{\mathcal A}(A,U\pitchfork B)
  =\operatorname{Hom}_k\bigl(U,
                \operatorname{Hom}_{\mathcal A}(A,B)\bigr).
\]
Let \([C,D]_{\ge0}\) be the usual connective internal Hom of
chain complexes: its positive-degree terms are those of the
full Hom complex, and its degree-zero term is the kernel of
the differential out of degree zero.  Tensor--Hom currying
identifies maps into this complex with chain maps from a tensor
product.  If \(D^1(U)=(U\xrightarrow{1}U)\) in degrees one
and zero, then a chain map \(E\to D^1(U)\) is determined
by an arbitrary map \(E_0\to U\); its degree-one component
is that map composed with \(d_E\).  Consequently an internal
Hom for (SM43.3) is
\[
 [(C,A),(D,B)]_{\mathcal M}=
 \left([C,D]_{\ge0}\oplus
          D^1\bigl(\operatorname{Hom}_{\mathcal A}(A,B)\bigr),
                         C_0\pitchfork B\right).             \tag{SM43.5}
\]
Indeed maps from \((E,F)\) into its three indicated summands
are exactly the chain map \(E\otimes C\to D\) and the two
maps \(E_0\odot A\to B\), \(C_0\odot F\to B\) specifying
a map \((E,F)\boxtimes(C,A)\to(D,B)\).

The chain category has enough projectives, generated by
\(k[0]\) and the projective disks \(D^n(k)\), \(n\ge1\).
Together with enough projectives in \(\mathcal A\),
componentwise projective epimorphisms prove that
\(\mathcal M\) has enough projectives.
The actual scalar copower of \(\mathbf1\) is
\[
 \kappa U=(U[0],0),\qquad
 \varphi(C,A)=H_0(C),\qquad \psi(C,A)=C_0,
                   \qquad\varphi\dashv\kappa\dashv\psi.
                                                               \tag{SM43.6}
\]
The two adjunctions follow by the ordinary factorization of
chain maps to degree zero through \(H_0\), and by evaluation
for maps from a degree-zero complex.  Maps involving the zero
second component contribute the one zero map.  The scalar
functor is exact, preserves all limits and colimits, and is
strong symmetric monoidal by (SM43.3).

Let \(\Lambda(C,A)=C\).  There is a natural canonical
identification
\[
                  L\varphi(C,A)\simeq C.                    \tag{SM43.7}
\]
To check it, resolve \((C,A)\) by projective objects
\((P_r,Q_r)\) of \(\mathcal M\).  Each \(P_r\) is
projective in the chain category, hence has zero positive
homology and projective \(H_0\).  The bicomplex \((P_r)_n\)
has a first-quadrant total complex.  Its resolution-direction
augmentation to \(C\) is a quasi-isomorphism, because it is
exact in that direction in every chain degree.  Its other
augmentation to \(H_0(P_\bullet)\) is a quasi-isomorphism,
because each \(P_r\) has zero positive homology.  Every total
degree is a finite sum, so this is precisely the finite-diagonal
argument of BR-13401, with no AB4 assumption on
\(\mathcal A\).  Applying \(\varphi\) to the projective
resolution gives \(H_0(P_\bullet)\), which proves (SM43.7)
with its natural comparisons.  On \(\kappa U\), the canonical
counit is consequently the identity of \(U[0]\), for every
vector space \(U\).

Take \(J=(D^1(k),0)\), and let
\(\alpha:\mathbf1\to J\) be the degree-zero inclusion.
Then
\[
 J\boxtimes(C,A)=(D^1(k)\otimes C,A),\qquad
 (J/\mathbf1)\boxtimes(C,A)=(k[1]\otimes C,0).              \tag{SM43.8}
\]
The map \(\alpha\boxtimes1\) is the identity in the second
component and degreewise monic in the first.  Write the latter
chain complex in degree \(n\) as \(C_n\oplus C_{n-1}\).
Its differential and a contraction are
\[
 d(x,y)=(d_Cx+y,-d_Cy),\qquad h(x,y)=(0,x),
                     \qquad dh+hd=1.                         \tag{SM43.9}
\]
Thus (SM43.7) gives \(L\varphi(J\boxtimes X)=0\) in
every degree, including zero, for every \(X\).  Moreover both
functors in (SM43.8) are exact.  This verifies the full interval
and even the optional flat-quotient feature, independently of
exactness of arbitrary coproducts in \(\mathcal A\).

We now choose the category \(\mathcal A=\mathcal N\) of
F-SM-34, whose abelian structure, all limits and colimits, and
enough projectives were proved there.  Recall the exact
countable non-AB4 witness.  With \(E(V)=V^{**}/V\), put
\[
 P_0(V)=(V,V,1,0),\qquad
 P_1(V)=(EV,V,0,1),\qquad V=k^{(\mathbb N)}.
\]
Since \(Ek=0\), each map
\(P_1(k)\to P_0(k)\) with components \((0,1)\) is monic.
Their countable coproduct is \(P_1(V)\to P_0(V)\), because
both \(P_i\) are free functors and preserve coproducts.
This map has kernel \((EV,0,0,0)\ne0\).  Nonvanishing follows
by pulling a nonzero linear functional on
\(k^{\mathbb N}/k^{(\mathbb N)}\) back to \(k^{\mathbb N}\):
it gives a bidual class which no evaluation vector represents.
Embedding these maps as \((0,P_1(k))\to(0,P_0(k))\) in
\(\mathcal M\) preserves this coproduct and its nonzero kernel.
Thus countable AB4 fails in \(\mathcal M\).

Finally, by (SM43.7), the weak equivalences are exactly the
maps whose first components are quasi-isomorphisms.
Projection \(\Lambda\) has the section \(G(C)=(C,0)\), with
\(\Lambda G=1\), and the natural map
\((C,A)\to(C,0)\) is a weak equivalence.  After localization
these maps yield mutually inverse functors, proving
\[
                     W^{-1}\mathcal M\simeq D_{\ge0}(k).
                                                               \tag{SM43.10}
\]
All coproducts in \(\mathcal M\) have the usual chain
coproduct as first component.  Hence (SM43.7) also identifies
each canonical map
\(\bigoplus_iL\varphi(X_i)\to
L\varphi(\bigoplus_iX_i)\) with the identity of that chain
coproduct.  Thus the derived coproduct condition is strictly
weaker than AB4 in this full interval setting.  This example
does not decide whether the original assumptions always force
that weaker condition or all-free scalar normalization.
\(\square\)

\paragraph{Source and dependence scope.}
The original predicates are from Grothendieck, A. (n.d.),
\emph{Pursuing stacks: Linearization of homotopy types},
cote~134-8, footers~46--48,
\url{https://grothendieck.umontpellier.fr/134-8.pdf}.
The abelian non-AB4 category is exactly the previously
fully proved F-SM-34; the derived identification uses the
already proved finite-diagonal lemma BR-13401.  The tensor,
all its coherence, closedness, the interval, and the
localization calculation are verified above.
